\section{Introduction}\label{sec:intro}

Problem 131 of Kagey's Open Problem Collection~\cite{kagey131} asks about
a Galton board where the ball tends to keep going the same way. Each bounce
repeats the one before it with probability $p$, and the question is when
the middle bin becomes as likely as an endpoint. The companion papers treat
this question for 2 directions~\cite{paperA} and for $d$ directions with
uniform switching~\cite{paperB}. In both, the crossing comes when $N$ times
the probability of switching to a given other direction is
$\log N-\frac12\log\log N+O(1)$, with an explicit constant in the $O(1)$
term. In this paper we ask the same question when the switching rule is a
general finite Markov chain.

\paragraph{The model.}
Let $Q$ be a generator on $[d]=\{1,\dots,d\}$, with rates $q_{ij}\ge0$ for
$i\ne j$ and exit rates $q_i=\sum_{j\ne i}q_{ij}$, and let $\mu$ be a
starting law on $[d]$. We run the chain $X_1,\dots,X_N$ with $X_1\sim\mu$
and transition matrix $I+hQ$, where
\[
 h=\frac TN,\qquad T=\tau\log N,\qquad \tau>0 .
\]
So at each step the chain leaves state $i$ with probability $hq_i$, and
$N$ steps look like the continuous-time chain with generator $Q$ run for
time $T$. The composition $n=(n_1,\dots,n_d)$ counts the visits to each
state, and it plays the role of the bin. For $d=2$, $q_{12}=q_{21}=1$ and a
uniform start, this is the walk of Paper~A with $h=1-p$. With $q_{ij}=1$ for
all $i\ne j$ and a uniform start it is the model of Paper~B
(Example~\ref{ex:models}). The
support of $n$ is a face $F$ of the simplex, and Kagey's question becomes
which face holds the most likely composition, and at which $T$ the mode
moves from one face to another. The middle bin against an endpoint is the
full face against a vertex.

\paragraph{The switching graph.}
Let $G_Q$ be the switching graph, with an edge $i\to j$ when $q_{ij}>0$.
For a face $F$ let $Q_F$ be the principal submatrix of $Q$ on $F$, and put
$\lambda_F=-\varrho(Q_F)$, where $\varrho$ is the largest real eigenvalue.
This is the exit rate of $F$. The continuous-time chain stays in $F$ up to
time $t$ with probability about $e^{-t\lambda_F}$. We show that if some
path of positive probability has support exactly $F$, then the most likely
composition with support $F$ has probability $N^{-\Phi_\tau(F)+o(1)}$, where
\[
 \Phi_\tau(F)=\tau\lambda_F+\lvert F\rvert-1 .
\]
The reason is that a path with support $F$ must enter each of the other
$\lvert F\rvert-1$ states of $F$, and each first entry is a switch of
probability about $h=T/N$. So entering the new states is a factor of about
$N^{-(\lvert F\rvert-1)}$, up to powers of $T$. A later run in a state needs
another switch,
but it also multiplies the number of ways to split the visits to that
state into runs by a factor of order $N$, so it is not penalized. Staying
inside $F$ for the whole time $T$ is a factor about $e^{-T\lambda_F}$. With
$T=\tau\log N$ both factors are powers of $N$. So a sticky walk prefers to
sit in a well-separated community, a set of states with a small exit rate
for its size. The winning faces are read off from the lower-left boundary
of the convex hull of the points $(\lambda_F,\lvert F\rvert-1)$, and the
winning face changes at
values of $\tau$ equal to minus the slopes of its edges
(Theorem~\ref{thm:hull} and Figure~\ref{fig:cycles}).

\paragraph{Results.}
Theorem~\ref{thm:first} proves this for every start $\mu$. Call a
face admissible if some path of positive probability has support exactly
$F$. With $L=\log N$ we prove
\[
 \log\max_nP_N(n)=-L\min_F\Phi_\tau(F)+O(\log L),
\]
where the minimum is over admissible faces, and that for large $N$ every
global mode has its support on a minimizing face. Note that $\Phi_\tau$
does not depend on $\mu$, so to first order the start enters only through
the set of admissible faces. If the minimizing faces are irreducible, meaning that
their switching graphs are strongly connected, the mode lies within
$O(1/L)$ of the quasi-ergodic law of the chain killed on leaving the face.
This holds for example when $Q$ has symmetric support or $\mu$ has full
support (Theorem~\ref{thm:hull}).

Section~\ref{sec:secondorder} finds the $\log\log N$ term. Say that $F$
satisfies (H) if it is irreducible and $\mu(F)>0$, write $k=\lvert F\rvert$,
and let $M_F$ be the largest probability of a composition with support $F$.
We prove a local law for the compositions in the interior of $F$
(Theorem~\ref{thm:local} and Lemma~\ref{lem:discrete}) and deduce that
\[
 M_F=N^{-(k-1)}\,T^{(k-1)/2}K_F\,e^{-T\lambda_F}\bigl(1+O(1/T+T^2/N)\bigr)
\]
uniformly for $1\le T\le N^{1/3}$, with an explicit constant $K_F$
(Theorem~\ref{thm:facemax}). The factor $T^{(k-1)/2}$ appears because the
occupation fractions fluctuate on the scale $T^{-1/2}$ in $k-1$ directions
inside the face. Now take 2 faces $F$ and $G$ with (H), and put
$\Delta=\lvert F\rvert-\lvert G\rvert\ge1$ and $\delta=\lambda_G-\lambda_F>0$.
Then every $T_N\in[1,N^{1/3}]$ with $M_F=M_G$ satisfies
\[
 \delta T_N+\frac\Delta2\log T_N=\Delta L-\log\frac{K_F}{K_G}
 +O\Bigl(\frac1L\Bigr)
\]
(Theorem~\ref{thm:crossing}). Hence
$T_N=\frac\Delta\delta\bigl[L-\frac12\log L+O(1)\bigr]$, and the
coefficient $\frac12$ of $\log\log N$ in Papers~A and~B holds for every such
pair, whatever the graph, the rates, the sizes or the start, and without
reversibility. When all the faces that tie to first order satisfy (H), the
global mode follows the upper envelope of explicit lines in a window of
this form (Theorem~\ref{thm:window}). Papers~A and~B are special cases
(Corollaries~\ref{cor:paperA} and~\ref{cor:paperB}). For reversible faces
$K_F$ involves the square root of a weighted count of spanning trees
(Proposition~\ref{prop:constant}), and for Paper~B a weighted form of
Cayley's formula evaluates it. Without (H) the coefficient can change. On a
directed 3-cycle the mode passes through a reducible face, and the 2
crossings have coefficients 1 and 0 (Proposition~\ref{prop:threecycle}).

Section~\ref{sec:examples} works out examples. On the complete graph every
face has the same first-order exponent at $\tau=1$, which is the
degeneracy of Paper~B, and the second order decides. There the start
matters. On $K_3$ with $\mu=(\frac12,\frac12,0)$ an edge is the mode on a
window of width about $0.026$ in the offset of Theorem~\ref{thm:window},
and with a uniform start this window is empty
(Proposition~\ref{prop:K3start}). On the cycle $C_d$ with $d\ge4$ the
winners are arcs of $1,2,\dots,\lfloor2d/3\rfloor$ states and then the
whole cycle (Theorem~\ref{thm:cycle}). Once an arc covers two thirds of the
cycle, closing the cycle beats growing the arc. Two triangles joined by a
link of small rate $\varepsilon$ form 2 communities, and a triangle wins
from $\tau\approx1$ to $\tau\approx9/\varepsilon$
(Proposition~\ref{prop:triangles}).

Section~\ref{sec:sticky} applies face selection to sticky priors for hidden
Markov models. We fix the exit probability of every state at
$h=c\log N/N$ and let the
switching kernel $\Pi$ be a stochastic matrix with zero diagonal and
positive off-diagonal entries. Then for large $N$ the most likely
composition under the prior, which we call the MAP composition, has its
support on a minimizing face, with no reversibility needed
(Corollary~\ref{cor:map}). For $d\ge3$ the minimizing faces go from single
states straight to all $d$ states exactly when
$\varrho(\Pi_F)\le(\lvert F\rvert-1)/(d-1)$ for every face $F$
(Proposition~\ref{prop:thresholds}). For a symmetric kernel this happens
only for the uniform kernel of Paper~B. Under the weak-limit sticky HDP-HMM
prior with 3 states and uniform base weights, it has probability
$\frac{3\sqrt3}{8\pi\omega}(1+o(1))$ as the concentration $\omega\to\infty$
(Proposition~\ref{prop:directjump}). So the degeneracy of Paper~B is not
typical. For a symmetric pair of states with equal start weights, a single
state against the pair is exactly Kagey's crossing
(Proposition~\ref{prop:sticky-pair}). The MAP composition can also mislead.
Below the first threshold it uses a single state, while the event that the
chain uses a single state has probability about $N^{-c}$
(Proposition~\ref{prop:mass}).

Section~\ref{sec:planar} treats the planar persistent walk, which is
Kagey's Galton board in 2 dimensions. The walk moves in the directions
$E,N,W,S$. At each step it keeps its direction with probability
$1-(2r+s)h$, turns left or right with probability $rh$ each, and reverses
with probability $sh$. For even $N$ we compare the corner $(N,0)$, which is
reached only by going east at every step, with the origin. The crossing is
unique (Proposition~\ref{prop:decomp}). For $r>0$ it happens at
$h=T^*_N/N$ with $T^*_N/\log N\to\min\bigl(2/(2r+s),1/s\bigr)$
(Theorem~\ref{thm:planar-first}). The first value comes from paths that use
all 4 directions and the second from paths along a single axis, so the
regime switches at $s=2r$. When reversals dominate ($s>2r$), $sT^*_N$
follows the expansion of Paper~A with $\log N$ replaced by $\log N-\log2$,
since the 2 axes double the mass at the origin
(Theorem~\ref{thm:planar-rev}). When turns dominate ($s<2r$), a
two-dimensional local limit (Theorem~\ref{thm:fourdir}) gives
\[
 (2r+s)T^*_N+\log T^*_N=2\log N+\log\frac{\pi}{8(r+s)}+o(1)
\]
(Corollary~\ref{cor:planar-turn}). With turns only ($s=0$), $rT^*_N$ agrees
with the crossing $N(1-p_N)$ of Paper~A up to $O((\log N)^4/N)$
(Proposition~\ref{prop:pureturns}).

Section~\ref{sec:fisher} asks how much the endpoint alone tells us about the
switching rate. For the walk of Paper~A with flip probability $q$, let
$e_N(q)$ be the Fisher information about $q$ in the endpoint divided by the
Fisher information in the whole path. With $x=Nq$ we prove
$e_N(q)=\frac1{2x}+\frac1{4x^2}+O(x^{-3}+x^3/N)$ for $1\le x\le N^{1/4}$
(Theorem~\ref{thm:fisher}). At Kagey's crossing this gives
$e_N\approx1/(2\log N)$ (Corollary~\ref{cor:fisher-crossing}), but the
correction is large. At $N=3200$ the product $2e_N\log N$ is still about
$1.35$.

\paragraph{Relation to established results.}
The rate $I_F$ of Section~\ref{sec:firstorder} is the Donsker--Varadhan
rate of the occupation measure of the chain killed on leaving $F$. For
reversible jump processes Dupuis and Liu give the Donsker--Varadhan rate
explicitly~\cite[Eq.~(3.2)]{dupuisliu2015}. On a finite reversible chain it
reads $\sum_i\alpha_iq_i-\sum_{i\ne j}\sqrt{\alpha_i\alpha_jq_{ij}q_{ji}}$ for
a law $\alpha$ on the states, and for $\alpha$ carried by $F$ this is
$I_F(\alpha)$. Guillin, Nectoux and Wu~\cite{gnw2024} prove a large
deviation principle for the occupation measure of a process conditioned to
stay in a set (their Thm~2), and its rate vanishes only at the
quasi-ergodic law (their (2.13)). When the process is reversible they prove
that the exit rate of the set is the least value of the Donsker--Varadhan
rate over laws on the set (their Cor.~2). In general they prove an
inequality between their rate, the Donsker--Varadhan rate and the exit rate
(their (4.6)), and they conjecture equality in (4.6) (their (2.15)). For a
finite chain their condition (C5) requires $Q_F$ to be irreducible. For such
a chain the exit rate is the least value of $I_F$ also without
reversibility. We give a short self-contained proof in
Proposition~\ref{prop:rate}(a), because we also need the other properties of
$I_F$ proved there. It does not address their conjecture (2.15), which
concerns their conditional rate. The quasi-ergodic law goes back to Darroch and
Seneta~\cite{darrochseneta1965,darrochseneta1967} for irreducible absorbing
chains in discrete and continuous time, and Colonius and
Rasmussen~\cite{coloniusrasmussen2021} treat the reducible case.

Dietz and Sethuraman~\cite{dietz2007} study chains that switch from $i$ to
$j$ at time $n$ with probability $G(i,j)/n$, where every $G(i,j)>0$. Their
occupation fractions converge in law to a distribution with full support on
the simplex, and it is a Dirichlet law when $G(i,j)$ depends only on $j$.
Their chains also make about $\log n$ switches by time $n$, but the
switches are spread log-uniformly in time and the switching graph is
complete. In an earlier paper~\cite{dietz2005} they prove large deviations
at speed $n$ for chains whose transition matrices converge to a reducible
matrix, and they remark that scales up to $\log n$ are of a different kind.

Exact distributions of transition counts go back to
Whittle~\cite{whittle1955}, Dawson and Good~\cite{dawsongood1957} and
Goodman~\cite{goodman1958}. Polettini~\cite{polettini2015} and Carugno,
Vivo and Coghi~\cite{cvc2022} derive such counts from the BEST theorem and
the matrix-tree theorem and study the statistics of transition counts when
every count grows linearly with the length of the path. We group the words
by runs instead (Lemma~\ref{lem:runs}), which suits a regime with only
$O(\log N)$ switches. Bertini, Gabrielli and Landim~\cite{bgl2024}, and
Landim~\cite{landim2023} for non-reversible chains, expand the level-two
large deviation rate functions of finite chains whose rates depend on a
parameter. They describe the time scales of such chains but do not ask
where the mode of the composition sits. The sticky HDP-HMM of Fox,
Sudderth, Jordan and Willsky~\cite{fox2011} adds a self-transition bias to
the hierarchical Dirichlet process hidden Markov model, and
Section~\ref{sec:sticky} studies an idealized form of its prior.

On the lattice the correlated random walk goes back to
Gillis~\cite{gillis1955}. Chen and Renshaw study the Gillis--Domb--Fisher
walk and its $n$-step characteristic function~\cite{chenrenshaw1992}, and
they show that the characteristic function of any correlated random walk
satisfies a recurrence~\cite{chenrenshaw1994}, so in principle the law of
our planar walk is known. Marris and Giuggioli~\cite{marris2024} give
lattice propagators for walks with forward, backward and lateral steps and
solve their first-passage problem. In the continuum, Santra, Basu and
Sabhapandit~\cite{santra2020} find the exact marginals of a run-and-tumble
particle with 4 directions that turns by $\pm\pi/2$ and never reverses,
and they describe from simulations how its mass moves from the boundary to
the center. Angelani~\cite{angelani2022} treats general switching among 4
orthogonal directions, and his walk with turns only factorizes in
coordinates rotated by $45^\circ$~\cite[Eq.~(36)]{angelani2022}.
Proposition~\ref{prop:pureturns} uses a lattice analogue of this
factorization. In the continuum the corners carry atoms and the origin only
a density, so the comparison of a corner with the origin is a lattice
question. Smiga, Radaelli, Binder and Landi~\cite{smiga2023} compute the
Fisher information of the empirical distribution of a stationary process
under a Gaussian approximation. Their leading term depends on how the
stationary law moves with the parameter. For the symmetric chain of
Paper~A the stationary law does not depend on $q$, so their leading term
vanishes here. We did not find the face-selection question or the planar
crossing of the corner and the origin in the literature.

\paragraph{Organization.}
Section~\ref{sec:model} sets up the model and the exact law.
Section~\ref{sec:firstorder} proves first-order face selection and
describes the winners, and Section~\ref{sec:secondorder} gives the local
law, the face maxima and the crossing equation.
Sections~\ref{sec:examples}, \ref{sec:sticky}, \ref{sec:planar}
and~\ref{sec:fisher} treat the examples, the sticky priors, the planar walk
and the Fisher information. Section~\ref{sec:verify} describes the
numerical checks and lists open questions, and the longer proofs are in the
appendices. We do not treat compositions near the boundary of a face,
where a local law that holds uniformly is open (Example~\ref{ex:star}), and
we do not prove in general that a crossing is unique.
